\documentclass[12pt]{letter}
\usepackage[margin=2.5cm]{geometry}
\usepackage{hyperref}
\usepackage{url}
\urlstyle{same}

\signature{Firman Hadi\\
Department of Geodetic Engineering\\
Universitas Diponegoro\\
Semarang, Indonesia\\
\url{firmanhadi21@lecturer.undip.ac.id}}

\address{Firman Hadi\\
Department of Geodetic Engineering\\
Faculty of Engineering\\
Universitas Diponegoro\\
Semarang 50275, Indonesia\\
\url{firmanhadi21@lecturer.undip.ac.id}}

\date{March 20, 2026}

\begin{document}

\begin{letter}{%
Editors\\
\emph{Transportation Research Interdisciplinary Perspectives}}

\opening{Dear Editors,}

We are pleased to submit our manuscript, ``Spatiotemporal Analysis of Urban Traffic Congestion,
Network Topology, and Capacity Constraints in Indonesian Metropolitan Cities,'' for
consideration in \emph{Transportation Research Interdisciplinary Perspectives}.

\textbf{Relevance to the journal.}
This manuscript bridges spatial analysis, network science, traffic engineering, and urban
planning---making it a natural fit for TRIP's interdisciplinary scope. Our central finding---that
congestion operates as a demand synchronization phenomenon rather than a spatial infrastructure
problem---challenges conventional supply-side responses and carries important policy implications
for rapidly urbanising cities in the Global South. We believe this aligns with TRIP's stated
openness to results that question established assumptions.

\textbf{Summary.}
We analyze 264 million traffic observations collected at 15-minute intervals over 11 months
across three Indonesian metropolitan cities---Jakarta (metro pop.\ 34M), Bandung (2.5M), and
Semarang (1.8M)---spanning a 20$\times$ population range. Using multilevel mixed-effects models
on absolute speed, combined with geostatistical methods (Moran's I, LISA, Getis-Ord Gi*),
network topology analysis, and graph-based capacity drop detection, we empirically decompose
congestion into its temporal and spatial components.

Time-of-day explains 57--67\% of within-segment speed fluctuations, while network centrality
adds less than 1\% beyond road type. A positive-control test confirms that the same analytical
pipeline detects spatial structure in road design ($R^2 = 0.60$--$0.71$ for free-flow speed)
but not in congestion ($R^2 < 0.003$), ruling out methodological artefacts. Cross-metric
validation using absolute speed, jam factor, and speed reduction confirms these patterns are
not normalization artefacts. H3 hexagonal aggregation at multiple scales further confirms
the null spatial findings are not modifiable areal unit artefacts.

\textbf{Contributions.}
\begin{enumerate}
\item Empirical evidence that congestion is a demand synchronization phenomenon, with temporal
factors explaining 57--67\% of speed variation while network topology adds less than 1\%.
\item A positive-control methodology demonstrating that the analytical pipeline detects spatial
structure in road design but not in congestion, ruling out methodological artefacts.
\item Cross-metric validation showing results hold for absolute speed, addressing the circularity
inherent in capacity-normalized congestion metrics.
\item A fully open, installable Python pipeline (\texttt{pip install traffic-congestion-pipeline})
bridging commercial traffic data with the OSMnx/PySAL ecosystem.
\end{enumerate}

\textbf{Open science resources:}
\begin{itemize}
\item \textbf{Source code:} \url{https://github.com/firmanhadi21/traffic-analyses} (MIT licence)
\item \textbf{PyPI package:} \url{https://pypi.org/project/traffic-congestion-pipeline/}
\item \textbf{API documentation:} \url{https://firmanhadi21.github.io/traffic-analyses/}
\item \textbf{Archived datasets:} Zenodo DOI: \href{https://doi.org/10.5281/zenodo.18650759}{10.5281/zenodo.18650759}
\end{itemize}

\textbf{Confirmations.}
This manuscript is original and has not been published or submitted elsewhere. All authors have
approved the manuscript and agree with submission to \emph{Transportation Research
Interdisciplinary Perspectives}. We have no competing interests to declare. The use of
generative AI (Claude, Anthropic) for code development and statistical automation is disclosed
in the manuscript per Elsevier policy.

\closing{Sincerely,}

\end{letter}
\end{document}
